\chapter{Conclusion (Provisional)}
\label{ch:conclusion}

\todo[inline]{Provisional conclusion: rewrite once M4--M7 are done and the research question can be answered with repeated runs and confidence intervals.}

This thesis builds a graph-specific, cost-aware agent harness and asks whether it lets a cheap model match or beat a
frontier model in a general-purpose harness at a fraction of the cost. On the date of this draft the harness itself is
built: a loop with budgets and robust handling of malformed replies, validated tools that return evidence, structured
answers, verifiers that check evidence without becoming oracles, result handles and bounded messages that keep the
context independent of the graph size, an MCP server that gives other harnesses exactly the same tools, a locked-down
sandbox for model-written code, context compaction, and logging and analysis from which every number can be
reproduced.

The preliminary results are consistent with the hypothesis on the cost side and inconclusive on the accuracy side.
On the existing dataset, a local 8-billion-parameter model in the harness answered everything correctly in lenient
mode, and its remaining errors were format errors that the harness can rescue or, in future, prevent. A
general-purpose harness with a frontier model reached the same accuracy on a small preview with 17--35 times more input
tokens per question. Well-fitting tools reduced cost by about an order of magnitude at equal accuracy, and on tasks
that need many lookups, code over the same tools turned near-total failure into near-total success for a cheap model.
These results come from single runs on few questions of the development split, mostly with one model, on a dataset that
is saturated on its largest graphs; they are not yet evidence for the research question.

What remains is what will decide it: harder and larger tasks from the scaled generator and the external benchmarks,
the same-model and same-tools comparison with general-purpose harnesses, the router, skills and escalation, and the
repeated runs with confidence intervals that let the accuracy-versus-cost Pareto front be drawn.
